\section{Existující řešení}
Následující podkapitola je věnována popisu existujících řešení, která umožňují sledovat a případně i porovnávat výkonnost vývojářů. Mezi tato řešení patří i samotné nástroje GitLab a GitHub, které nabízejí bezplatné zobrazení základních statistik o hostovaných repozitářích. Kromě těchto dvou nástrojů budou v následujících sekcích představena i řešení z komerční sféry, která mi přišla nejrozšířenější nebo nějakým způsobem zajímavá.

\subsection{GitLab}
Nástroj GitLab nabízí několik vizualizací, díky nimž je možné získat přehled o prováděných aktivitách na projektové, skupinové a instanční úrovni. Projektová úroveň odpovídá jednomu konkrétnímu repozitáři, skupinová zahrnuje jeden nebo více souvisejících projektů a instační obsahuje všechny projekty v daném systému.

V bezplatné variantě jsou avšak k dispozici pouze základní statistiky. Jedná se například o informace o využívaných programovacích jazycích, o počtu provedených commitů nebo o době trvání a úspěšnosti doběhnutí pipeline. \cite{GitLab-analytics} Zajímavější statistiky jsou k dispozici u placených plánů, jejichž cena začíná na 29 dolarech za uživatele na měsíc. \cite{GitLab-pricing} Ty následně zahrnují například statistiky o merge requestech a code review. U merge requestů se jedná o (průměrný) čas do mergnutí, počet mergnutých merge requestů za zvolené časové období a například počet commitů a změněných řádků u jednotlivých merge requestů. \cite{GitLab-merge-requests} Statistiky code review obsahují mimo jiné délku review procesu (čas od přidání prvního komentáře jinou osobou než autorem merge requestu), počet komentářů nebo seznam approvers. \cite{GitLab-code-review}

\subsection{GitHub}
TODO

\subsection{Swarmia}
https://www.swarmia.com/
